% Part: first-order-logic
% Chapter: beyond
% Section: other-logics

\documentclass[../../../include/open-logic-section]{subfiles}

\begin{document}

\olfileid{fol}{byd}{oth}

\olsection{ఇతర తర్కాలు}

ఇప్పటికే మీరు గ్రహించి ఉండవచ్చు: కొత్త తర్కాన్ని రూపొందించడం కష్టం
కాదు. మీరు కూడా మీ సొంత వాక్యనిర్మాణాన్ని సృష్టించి, ఒక నిగమన
వ్యవస్థను కల్పించి, దానికి సరిపోయే అర్థవిచారాన్ని రూపొందించవచ్చు.
ఆ అర్థవిచారానికి \tetoken{వ్యుత్పత్తి}{!!{derivation}} వ్యవస్థ సంపూర్ణంగా ఉండాలంటే కొంత
చాకచక్యం అవసరం కావచ్చు; మీ తర్కం నిజంగా ఆసక్తికరమైనదని ప్రపంచాన్ని
ఒప్పించడానికి కొంత శ్రమ పడవచ్చు. కానీ దానికి ప్రతిఫలంగా, మీ తర్కపు
గణిత, గణనాత్మక ధర్మాలను అన్వేషిస్తూ గంటలకొద్దీ చక్కని వినోదాన్ని
పొందవచ్చు.

ఇటీవలి దశాబ్దాల్లో ఔపచారిక తర్కాలు నిజంగా విస్ఫోటకంగా విస్తరించాయి.
అస్పష్ట ధర్మాల గురించిన హేతుచింతనను నమూనీకరించడానికి ఫజీ తర్కాన్ని
రూపొందించారు. అనిశ్చితి గురించిన హేతుచింతనను నమూనీకరించడానికి
సంభావ్యతా తర్కాన్ని రూపొందించారు. కొత్త సమాచారం ఎదురైనప్పుడు
తరువాత రద్దు చేయగల ``సమంజసమైన'' నిగమనాలు, అంటే రద్దుచేయదగిన హేతుచింతన
రూపాలను నమూనీకరించడానికి డిఫాల్ట్ తర్కాలు, ఏకదిశేతర తర్కాలను
రూపొందించారు. జ్ఞానం గురించిన హేతుచింతనను నమూనీకరించే
జ్ఞానమీమాంసాత్మక తర్కాలు; కారణ సంబంధాల గురించిన హేతుచింతనను
నమూనీకరించే కారణ తర్కాలు; నైతిక బాధ్యతల గురించిన హేతుచింతనను
నమూనీకరించే ``కర్తవ్యాత్మక'' (డియాంటిక్) తర్కాలు కూడా ఉన్నాయి. ఈ
వ్యవస్థలను ప్రవేశపెట్టడానికి ప్రధాన ప్రేరణ తాత్త్వికమా, గణితపరమా,
గణనాత్మకమా అనేదాన్ని బట్టి, ఇటువంటి వ్యవస్థలను గణిత తర్కం, తాత్త్విక
తర్కం, కృత్రిమ మేధ, సంజ్ఞాన శాస్త్రం లేదా మరేదైనా విభాగం కింద
అధ్యయనం చేయడాన్ని చూడవచ్చు.

ఈ జాబితా కొనసాగుతూనే ఉంటుంది; సాధ్యతలు అంతులేనివిగా కనిపిస్తాయి.
మానవ హేతుచింతన మొత్తాన్ని గణనకు తగ్గించాలనే లైబ్నిజ్ కలను మనం
ఎప్పటికీ సాధించకపోవచ్చు---కానీ ప్రయత్నించకుండా అది మనల్ని ఆపలేదు.

\end{document}
